% Requires amsmath, amssymb, amsthm and standard theorem environments.
% References to add: repolattice, okudaueno, dlmfbeta,
% dlmfgamma.  Bibitems and precise source coverage are in
% two_cluster_source_status.txt.  Every label has the prefix tc:.

\section{An explicit polylogarithm dictionary for two mass clusters}
\label{tc:sec:clusters}

The incoming \emph{Lattice Contacts and Mixed Zeta Identities} report
evaluates a Gamma face when all but one mass agree.  Its spectral-face
section then asks for a finite dictionary for two distinct mass
clusters; its research question~Q4 asks more broadly about displaced
masses and the dependence of depth on the number of clusters
\cite{repolattice}.  We evaluate the entire two-cluster sector,
with arbitrary positive cluster multiplicities.  The result is a
finite combination of ordinary one-variable multiple polylogarithms
inside an absolutely convergent spectral sum.  In particular, two
equal displaced masses require depth at most three, independently
of the total number of factors.

The argument uses classical beta integration, aggregation of Dirichlet
parameters, the Gamma product, and coefficient extraction.  The
transformation under exchange of clusters is a specialization of
classical multiple-polylogarithm Landen formulas
\cite{okudaueno}; it is not asserted to be a new general
transformation law.  The contribution relative to the inspected
incoming source is the explicit all-multiplicity dictionary and its
convergence and depth statement.  Minimality of the resulting depth
and a bound depending only on the number of distinct clusters remain
unproved.

\subsection{The ordinary face and its independent regulators}

Let $k\ge2$, $b_1,\ldots,b_k\in\mathbb R$, and
\[
 \Sigma_k=\{(t_1,\ldots,t_k):t_i\ge0,\ \sum_i t_i=1\}.
\]
The measure $dt$ means $dt_1\cdots dt_{k-1}$ after eliminating
$t_k$; it has no additional Euclidean surface factor.  For nonempty
$T\subseteq\{1,\ldots,k\}$ set
\[
 t_T=\sum_{i\in T}t_i,\qquad
 b_T=\frac{\sum_{i\in T}t_i b_i}{t_T},
\]
and give $t_TF(b_T)$ the value zero when $t_T=0$.  Define
\begin{align}
 \mathcal H_k(F;t)
  &=\sum_{\varnothing\ne T\subseteq\{1,\ldots,k\}}
       (-1)^{k-|T|}t_TF(b_T),\label{tc:eq:face-integrand}\\
 I_k(F;b_1,\ldots,b_k)
  &=\int_{\Sigma_k}\mathcal H_k(F;t)
                    \frac{dt}{t_1\cdots t_k}.
                    \label{tc:eq:ordinary-face}
\end{align}
Constants and affine functions are annihilated.

These are ordinary integrals.  To see convergence for $F\in C^1$
on the mass interval, each perspective is Lipschitz because
\[
 \partial_{t_i}\{t_TF(b_T)\}
    =F(b_T)+F'(b_T)(b_i-b_T),\qquad i\in T.
\]
The alternating sum vanishes on every coordinate face, so
$|\mathcal H_k(F;t)|\le C\min_i t_i$.
On the region where $t_j$ is maximal, $t_j\ge1/k$ and
\[
 \min_i t_i\le\left(\prod_{i\ne j}t_i\right)^{1/(k-1)}.
\]
Consequently the integral is absolutely convergent.  The same argument
allows sufficiently small independent complex powers $t_i^{\alpha_i}$
and any fixed powers of their logarithms.

For $\alpha_i>0$, $A=\sum_i\alpha_i$, let
\[
 D_F(\boldsymbol\alpha;\boldsymbol b)
   =\frac{\Gamma(A)}{\prod_i\Gamma(\alpha_i)}
      \int_{\Sigma_k}F\!\left(\sum_i t_i b_i\right)
                            \prod_i t_i^{\alpha_i-1}\,dt.
\]
If $A_T=\sum_{i\in T}\alpha_i$, the classical beta integral
\cite{dlmfbeta} gives
\begin{equation}\label{tc:eq:independent-regulators}
 \frac{\Gamma(1+A)}{\prod_i\Gamma(1+\alpha_i)}
 \int_{\Sigma_k}\mathcal H_k(F;t)
                     \prod_i t_i^{\alpha_i-1}\,dt
 =\frac{1}{\prod_i\alpha_i}
   \sum_{\varnothing\ne T}(-1)^{k-|T|}
                  A_TD_F(\boldsymbol\alpha_T;\boldsymbol b_T).
\end{equation}
Here the notation on the right denotes the Dirichlet mean formed
only from the labels in $T$, not the random barycenter on the left.
The identity follows from
$\mathbb E[t_TF(b_T)]=(A_T/A)D_F(\boldsymbol\alpha_T;\boldsymbol b_T)$.
It is important that the $\alpha_i$ are independent until the
squarefree coefficient is taken.  The left side is holomorphic near
zero and has value $I_k(F;\boldsymbol b)$ there.

\subsection{One scalar beta integral for arbitrary multiplicities}

Let $r,s\ge1$, $k=r+s$, and write $B^r,C^s$ for $r$ copies of
$B$ followed by $s$ copies of $C$.  Put $\delta=C-B$ and
\begin{align}
 E_F(t)&=F(B+\delta t)-(1-t)F(B)-tF(C),\label{tc:eq:chord}\\
 Q(X,Y)&=\frac{\Gamma(1+X+Y)}{\Gamma(1+X)\Gamma(1+Y)}.
                                                     \label{tc:eq:Q}
\end{align}

\begin{theorem}[Scalar beta representation for two clusters]
\label{tc:thm:beta}
For $F\in C^2$ on the interval with endpoints $B,C$,
\begin{equation}\label{tc:eq:beta-representation}
 \boxed{\displaystyle
 I_{r+s}(F;B^r,C^s)
   =r!s!\,[X^{r-1}Y^{s-1}]\,
      Q(X,Y)\int_0^1
      \frac{E_F(t)}{t(1-t)}(1-t)^X t^Y\,dt.}
\end{equation}
The integral on the right is holomorphic, for example for
$\Re X,\Re Y>-1$, and the coefficient may be taken under the integral.
\end{theorem}

\begin{proof}
Aggregate the independent regulators belonging to the $B$ cluster
into $X$ and those belonging to the $C$ cluster into $Y$.  The
proportion of the $C$ cluster has beta density
$t^{Y-1}(1-t)^{X-1}/\mathrm B(Y,X)$.  Thus the full-set term
$AD_F$ from \eqref{tc:eq:independent-regulators} becomes
\begin{align}
 \mathcal M_F(X,Y)
 &=\frac{\Gamma(1+X+Y)}{\Gamma(X)\Gamma(Y)}
     \int_0^1 F(B+\delta t)t^{Y-1}(1-t)^{X-1}\,dt\notag\\
 &=XF(B)+YF(C)
   +XYQ(X,Y)\int_0^1
          \frac{E_F(t)}{t(1-t)}(1-t)^Xt^Y\,dt.
                                                     \label{tc:eq:M}
\end{align}
The second line follows by subtracting the endpoint chord and
integrating it exactly.  Since $E_F(t)=O(t(1-t))$, it proves
holomorphy of the full-set term near $(0,0)$ without leaving
unremoved beta poles.

In the alternating sum in \eqref{tc:eq:independent-regulators},
each proper-subset term is holomorphic by the same chord
subtraction, or is linear when it contains only one cluster.
Every proper subset omits at least one regulator.  Its coefficient
of $\alpha_1\cdots\alpha_k$ is therefore zero.  In the full-set term,
the corresponding coefficient is
$r!s![X^rY^s]\mathcal M_F(X,Y)$: the squarefree coefficient in the
$r$th power of a sum of $r$ independent variables is $r!$.
The chord term in \eqref{tc:eq:M} is linear and contributes nothing.
Taking the removable value of
\eqref{tc:eq:independent-regulators} proves
\eqref{tc:eq:beta-representation}.  Integrable powers at the two
endpoints justify every fixed coefficient extraction.
\end{proof}

All kernels in this formula are explicit finite polynomials.  Indeed,
\begin{equation}\label{tc:eq:logQ}
 \log Q(X,Y)=\sum_{n=2}^{\infty}
  \frac{(-1)^n\zeta(n)}n
      \bigl((X+Y)^n-X^n-Y^n\bigr).
\end{equation}
For instance, with $L=\log(1-t)$ and $M=\log t$,
\begin{align}
 [XY]Qe^{XL+YM}&=LM+\zeta(2),\label{tc:eq:kernel22}\\
 [X^2Y]Qe^{XL+YM}&=\tfrac12L^2M+\zeta(2)L-\zeta(3),\\
 [X^2Y^2]Qe^{XL+YM}
  &=\tfrac14L^2M^2+\zeta(2)LM-\zeta(3)(L+M)
      +\tfrac32\zeta(4)+\tfrac12\zeta(2)^2.
\end{align}
These provide ordinary log-Gamma integral identities without any
endpoint finite-part prescription.

\subsection{The finite composition sum}

Our one-variable multiple polylogarithm convention is
\begin{equation}\label{tc:eq:MPL}
 \operatorname{Li}_{a_0,a_1,\ldots,a_d}(z)
  =\sum_{n_0>n_1>\cdots>n_d\ge1}
      \frac{z^{n_0}}{n_0^{a_0}n_1^{a_1}\cdots n_d^{a_d}},
 \qquad |z|<1.
\end{equation}
Only the outer index receives $z$.  Its branch on
$\mathbb C\setminus[1,\infty)$ is continued from zero by the
usual iterated integrals; there is no branch singularity at zero.
For $r,s\ge1$, $k=r+s$, define the finite combination
\begin{equation}\label{tc:eq:P}
 \boxed{\displaystyle
 P_{r,s}(z)=\sum_{d=1}^{s}\binom{k-d-1}{r-1}
      \sum_{\substack{p_1+\cdots+p_d=k-1\\p_i\ge1}}
                \operatorname{Li}_{1,p_1,\ldots,p_d}(z).}
\end{equation}
Every term has weight $k$ and depth at most $s+1$.

\begin{lemma}[Exact logarithmic face]
\label{tc:lem:logarithm}
If $x>0$ and $B,C\ge0$, set
$q=(C-B)/(x^2+C)$.  Then
\begin{equation}\label{tc:eq:logarithmic-face}
 I_{r+s}\bigl(U\mapsto-\log(x^2+U);B^r,C^s\bigr)
     =(-1)^{r+s-1}r!s!P_{r,s}(q).
\end{equation}
\end{lemma}

\begin{proof}
First suppose $C\ge B$, so $0\le q<1$.  The logarithm is
\[
 -\log(x^2+C)-\log\bigl(1-q(1-t)\bigr).
\]
The constant contributes only a linear term to
$\mathcal M_F(X,Y)$.  Applying the beta moment to the convergent
logarithmic series gives
\begin{equation}\label{tc:eq:log-generating-function}
 \mathcal M_F(X,Y)+(X+Y)\log(x^2+C)
   =X\sum_{m=1}^{\infty}\frac{q^m}{m}
            \prod_{j=1}^{m-1}\frac{1+X/j}{1+(X+Y)/j}.
\end{equation}
This series and its parameter derivatives converge normally for
$|q|<1$ and $X,Y$ in a sufficiently small fixed neighborhood of
zero.  One may use the Gamma-ratio expression for the product to
obtain a polynomial bound in $m$, uniformly on compact parameter
sets; the geometric factor $q^m$ then dominates every derivative.

For a finite product,
\[
 \prod_{j=1}^{m-1}\frac{1+X/j}{1+(X+Y)/j}
   =\prod_{j=1}^{m-1}\left(1-\frac{Y}{j+X+Y}\right).
\]
Choose $d$ nonconstant factors and expand their denominators.
The coefficient of $X^{r-1}Y^s$ is
\begin{align}
 &(-1)^{k-1}\sum_{d=1}^{s}\binom{k-d-1}{r-1}
  \sum_{\substack{p_1+\cdots+p_d=k-1\\p_i\ge1}}
  \sum_{m>j_1>\cdots>j_d\ge1}
                \frac{1}{j_1^{p_1}\cdots j_d^{p_d}}.
                                                     \label{tc:eq:coefficient}
\end{align}
For clarity, the $d$ selected factors supply $(-Y)^d$.
Their denominator expansions supply a total power
$(X+Y)^{k-d-1}$ and sign $(-1)^{k-d-1}$.
The remaining coefficient is precisely
$\binom{k-d-1}{r-1}$.  This proves the sign and every binomial
coefficient in \eqref{tc:eq:coefficient}; no relation among zeta
values has been used.  Summing in $m$ identifies exactly
\eqref{tc:eq:P}.  The coefficient identity from
Theorem~\ref{tc:thm:beta} proves
\eqref{tc:eq:logarithmic-face} for $0\le q<1$.

Both sides are holomorphic in $q$ on
$\mathbb C\setminus[1,\infty)$.  On the integral side this follows
by subtracting the endpoint chord of $-\log(1-q(1-t))$;
on each compact subset the subtracted numerator is uniformly
$O(t(1-t))$.  The identity theorem therefore includes $q<0$,
which is the case $C<B$.  In particular, the analytic continuation
uses real logarithms and real multiple polylogarithms on the whole
real interval $q<1$.
\end{proof}

\subsection{Gamma faces, convergence, and depth}

For $a>0$ and $U\ge0$ define
\begin{equation}\label{tc:eq:G}
 G_a(U)=\log\Gamma(a+i\sqrt U)+\log\Gamma(a-i\sqrt U)
                      -\log(2\pi),
\end{equation}
where the two logarithms are continued from positive real arguments.
Their sum is real.  The classical Gamma product
\cite{dlmfgamma} gives
\begin{equation}\label{tc:eq:Gamma-product}
 G_a(U)=2\log\Gamma(a)-\log(2\pi)
          -\sum_{n=0}^{\infty}
                    \log\left(1+\frac{U}{(n+a)^2}\right).
\end{equation}

\begin{theorem}[Every two-cluster Gamma face]
\label{tc:thm:Gamma}
For $a>0$, $B,C\ge0$ and $r,s\ge1$,
\begin{equation}\label{tc:eq:Gamma-dictionary}
 \boxed{\displaystyle
 I_{r+s}(G_a;B^r,C^s)
   =(-1)^{r+s-1}r!s!
      \sum_{n=0}^{\infty}
       P_{r,s}\!\left(\frac{C-B}{(n+a)^2+C}\right).}
\end{equation}
The spectral sum is ordinary and absolutely convergent, locally
normally in the stated real parameters.  The same identity has a
dictionary of depth at most $\min(r,s)+1$, obtained if necessary by
exchanging the two clusters and using the continued branch on
the negative real axis.  This is an upper bound on a supplied
representation, not a claim of minimal functional depth.
\end{theorem}

\begin{proof}
Apply the ordinary face functional termwise to
\eqref{tc:eq:Gamma-product}.  This is justified directly in the
scalar beta formula: the endpoint-chord remainder of the $n$th
logarithm is bounded in $C^0$ by a constant times
$t(1-t)(C-B)^2(n+a)^{-4}$ on each compact parameter set.
The fixed logarithmic beta kernels are integrable.  Alternatively,
the second derivatives of the log terms converge normally and the
face is continuous in the $C^2$ norm.  Constants disappear, and
Lemma~\ref{tc:lem:logarithm} gives the displayed result.

Each $P_{r,s}(z)$ is $O(z^2)$ at zero, since its first index is one
and its depth is at least two.  The arguments in the tail are
$O(n^{-2})$.  Hence the spectral summands are $O(n^{-4})$,
uniformly on compact subsets.  A finite number of initial arguments
may lie to the left of $-1$ if $C<B$; their analytically continued
values are regular.  This proves absolute and local normal
convergence without a small-displacement hypothesis.

The first expression has depth at most $s+1$.
The same face with the two clusters exchanged has depth at most
$r+1$.  Taking the smaller establishes the asserted bound.
\end{proof}

The symmetry just used gives the precise transformation
\begin{equation}\label{tc:eq:Landen}
 P_{r,s}(z)=P_{s,r}\!\left(\frac{z}{z-1}\right),
       \qquad z\in\mathbb C\setminus[1,\infty).
\end{equation}
It also follows from the classical composition form of the
multiple-polylogarithm Landen connection formula
\cite{okudaueno}.  Here it provides a check on branch conventions
and a useful way to choose the smaller-depth expression.

The first collision coefficient checks all factorial normalizations:
\begin{equation}\label{tc:eq:collision-check}
 I_{r+s}(G_a;B^r,(B+\delta)^s)
  =\frac{(-1)^{r+s-1}}2rs(r+s-2)!\,
     \delta^2\sum_{n=0}^{\infty}
                     \bigl((n+a)^2+B\bigr)^{-2}
       +O(\delta^3).
\end{equation}
This agrees with the predecessor's factorial-variance monomial
formula.  The full dictionary is global for $B,C\ge0$, whereas
that local collision expansion need not converge at the desired
displacement.

\subsection{Two equal displaced masses and a worked ordinary integral}

For $s=1$, \eqref{tc:eq:P} becomes
$P_{r,1}=\operatorname{Li}_{1,r}$ and recovers precisely the
incoming one-displaced theorem.  For $s=2$ it gives the new
all-multiplicity specialization
\begin{equation}\label{tc:eq:two-displaced}
 \boxed{\displaystyle
 P_{r,2}(z)=r\operatorname{Li}_{1,r+1}(z)
     +\sum_{j=1}^{r}\operatorname{Li}_{1,j,r+1-j}(z).}
\end{equation}
Thus $r$ equal base masses and two equal displaced masses have
depth at most three for every $r$.  No increase in the base
multiplicity increases this bound.

For example, set
$E_a(t)=G_a(B+(C-B)t)-(1-t)G_a(B)-tG_a(C)$ and
$q_n=(C-B)/((n+a)^2+C)$.  The case $r=s=2$ gives the completely
ordinary identity
\begin{align}
 &\int_0^1\frac{E_a(t)}{t(1-t)}
                 \bigl(\log t\log(1-t)+\zeta(2)\bigr)\,dt
       \notag\\
 &\qquad=-\sum_{n=0}^{\infty}
   \left(2\operatorname{Li}_{1,3}(q_n)
         +\operatorname{Li}_{1,1,2}(q_n)
         +\operatorname{Li}_{1,2,1}(q_n)\right).
                                                     \label{tc:eq:worked-integral}
\end{align}
At $a=2.3$, $B=0.2$, $C=0.9$, the associated face
$I_4(G_a;B,B,C,C)$ is
\[
 -0.0836254300916287651503714290194933682217772138996833\ldots.
\]
This decimal is an independently checked diagnostic; the identity
is proved above for the full parameter domain.

\subsection{Mass derivatives and anchored polylogarithmic primitives}

Delete the first index one in \eqref{tc:eq:P} and put
\begin{equation}\label{tc:eq:Rpolylog}
 R_{r,s}(z)=\sum_{d=1}^{s}\binom{r+s-d-1}{r-1}
   \sum_{\substack{p_1+\cdots+p_d=r+s-1\\p_i\ge1}}
                       \operatorname{Li}_{p_1,\ldots,p_d}(z).
\end{equation}
The defining strict series gives
$P'_{r,s}(z)=R_{r,s}(z)/(1-z)$.  It follows that
\begin{align}
 \partial_C I_{r+s}(G_a;B^r,C^s)
   &=(-1)^{r+s-1}r!s!
       \sum_{n=0}^{\infty}
           \frac{R_{r,s}(q_n)}{(n+a)^2+C},\label{tc:eq:Cderivative}\\
 \partial_B I_{r+s}(G_a;B^r,C^s)
   &=-(-1)^{r+s-1}r!s!
       \sum_{n=0}^{\infty}
           \frac{R_{r,s}(q_n)}{(n+a)^2+B}.
                                                     \label{tc:eq:Bderivative}
\end{align}
These are total derivatives of the repeated cluster parameters:
$\partial_C$ moves all $s$ copies of $C$ together.  They should not
be confused with a derivative in just one independently labelled mass.
The displayed expressions have depth at most $s$; exchanging clusters
gives a representation with depth at most $\min(r,s)$.

Termwise differentiation is legitimate.  For the tail, the arguments
lie in one fixed small disk, where $P(z)=O(z^2)$ and $R(z)=O(z)$;
local holomorphic extensions and Cauchy estimates give normal
convergence of every fixed parameter derivative.  The finitely many
initial summands stay off their cuts.  Consequently these equations
also give higher polygamma identities after further differentiation.
Since $I_{r+s}(G_a;B^{r+s})=0$, integration of
\eqref{tc:eq:Cderivative} from $C=B$ to the desired endpoint is an
anchored antiderivative with no undetermined constant.

\subsection{Verification and the exact scope of the resolved question}

The supplied script compares two independently organized exact
coefficient calculations: the incoming monomial-face formula at
$0^r,1^s$, and \eqref{tc:eq:P} after the substitution
$z=u/(1+u)$.  It checks $r,s\in\{1,2,3,4\}$ and monomial degrees
$2$ through $15$, together with the corresponding Landen
coefficients, for a total of $448$ exact rational assertions.
Four numerical comparisons at $90$ decimal digits use the ordinary
beta integral on one side and strict harmonic sums with a
Hurwitz-zeta spectral tail on the other.  Their absolute residuals
range from $9.8\times10^{-77}$ to $3.3\times10^{-75}$.
One case has $C-B>a^2+B$, outside the predecessor's local
collision-series disk.  These are floating-point diagnostics,
not certified interval enclosures or proof-assistant proofs.

The result resolves the request at the end of the incoming
spectral-face section for an explicit dictionary when the mass set
has exactly two distinct values, allowing every multiplicity.  It
partially advances Q4 of that report.  It does not evaluate two
\emph{unequal} displaced masses above a third base mass, does not
prove that every two-cluster expression reduces to depth two, and
does not establish a minimal depth theorem.  A precise next target
is the three-cluster case $B^r,C,D$ with $C\ne D$, where an
aggregation into three independent regulator sums should retain
the same ordinary face normalization.
